Let $m\ge 1$ and $n\ge 2$.
An edge-coloring of $K_n$ using exactly $m$ colors is called \emph{balanced} if every vertex is incident to exactly $\lfloor n/m\rfloor$ edges of each color.

\textbf{Ambiguity / existence condition.} If $n\equiv 1\pmod m$ then $\lfloor n/m\rfloor=(n-1)/m=:q$ is an integer and the vertex-degree condition becomes ``each vertex has degree $q$ in each color class''. However this forces each color class to have exactly $nq/2$ edges, so a necessary condition for existence is that $nq$ is even (equivalently $m\mid \binom{n}{2}$). If $nq$ is odd, no balanced coloring exists and statements about ``every balanced coloring'' are vacuous.

\smallskip
Now fix a finite simple graph $G$ and let $m:=e(G)$.
Question: For which graphs $G$ is it true that for all sufficiently large $n$ with $m\mid (n-1)$ and for which balanced $m$-colorings of $K_n$ exist, \emph{every} balanced $m$-edge-coloring of $K_n$ contains a rainbow copy of $G$ (a copy of $G$ whose $m$ edges all have distinct colors)?
